% ---- Agent A section for the team Overleaf report (Yeseswini, BL.EN.U4AIE23013) ----
\subsection{Agent A: Viewpoint Planner (DQN)}
\label{sec:agentA}

\paragraph{Role.} Agent A decides where the UAV flies next and at which altitude, so that the whole bridge surface
is inspected with low detector uncertainty while spending as few steps and as little battery as possible.

\paragraph{MDP formulation.} The environment is an $8\times16$ grid of bridge-surface cells, 20\% of which contain a
hidden defect. The UAV flies at one of three altitudes (high, mid, low) with camera footprints of $5\times5$, $3\times3$
and $1\times1$ cells; lower altitude yields a sharper view and lower detector uncertainty. Each cell keeps its best view.
\begin{itemize}
  \item \textbf{State} $s_t \in [0,1]^{118}$: egocentric $7\times7$ windows of visited flags and uncertainty centred on the UAV,
        a $2\times4$ block summary of coverage and mean uncertainty, and the normalised row, column, altitude and battery.
        Since defects are hidden, the problem is a POMDP solved as an MDP over this observation.
  \item \textbf{Actions} $a_t \in \{\text{N},\text{S},\text{E},\text{W},\text{Climb},\text{Descend}\}$.
  \item \textbf{Reward} $r_t = 0.2\,\Delta U_t - 0.01 - 0.05\,\mathbb{1}[\text{no new information}] - \mathbb{1}[\text{out of bounds}]$,
        where $\Delta U_t$ is the total uncertainty removed (unseen cells count as 1). On success
        (coverage $\ge 95\%$ and mean uncertainty $\le 0.32$) the agent receives $10 + 20\,b_t$, with $b_t$ the remaining battery.
  \item \textbf{Termination}: success or empty battery; truncation after 200 steps.
\end{itemize}

\paragraph{Algorithm.} We use DQN~\cite{mnih2015human} with experience replay and a target network,
minimising $\big(r + \gamma \max_{a'} Q_{\bar\theta}(s',a') - Q_\theta(s,a)\big)^2$. Exploration is $\epsilon$-greedy with
$\epsilon$ decaying linearly from 1.0 to 0.05 over the first 30\% of training. Hyper-parameters: MLP $[256,256]$,
learning rate $5\times10^{-4}$, $\gamma=0.995$, buffer $10^5$, batch 128, target update every 1000 steps,
$2\times10^5$ environment steps (10.2 min on a Colab CPU). Implementation: Gymnasium and Stable-Baselines3~\cite{raffin2021sb3}.

\paragraph{Design iterations.} A first version with a global one-hot position (386 inputs) did not learn (0\% success);
an egocentric observation fixed this. A first reward with separate coverage and uncertainty terms allowed the agent to collect
reward without completing the inspection; rewarding total uncertainty reduction plus a battery-scaled success bonus aligned the
objective. A discount of $\gamma=0.99$ discounted the delayed success bonus too strongly, so we use $0.995$.

\paragraph{Results.} Table~\ref{tab:agentA} compares the learned policy with two passive raster scans on the same 20 seeded bridges.
DQN matches the dense low-altitude raster's 100\% success while using 29\% fewer steps and 26\% less energy. The trade-off is a
lower mIoU (0.539 vs.\ 0.626): the agent inspects just enough to satisfy the uncertainty criterion, which is precisely what the reward
specifies.

\begin{table}[h]
\centering
\caption{Agent A vs.\ raster baselines (analytic detector, 20 seeded episodes).}
\label{tab:agentA}
\begin{tabular}{lccc}
\hline
Metric & Raster mid & Raster low & \textbf{DQN} \\
\hline
Success rate & 0.00 & 1.00 & \textbf{1.00} \\
Steps & 199.0 & 123.0 & \textbf{87.5} \\
Energy used & 1.000 & 0.625 & \textbf{0.463} \\
Coverage & 1.000 & 0.953 & 0.998 \\
Defects found (/26) & 20.9 & 23.5 & 22.0 \\
mIoU & 0.500 & 0.626 & 0.539 \\
Return & 6.20 & 35.43 & \textbf{35.81} \\
\hline
\end{tabular}
\end{table}

\paragraph{Convergence and exploration.} Training return rises from $-25$ to about 34 and success reaches 100\% after roughly
750 episodes; most of the improvement happens after $\epsilon$ reaches its floor, i.e.\ once the replay buffer contains diverse
experience. Evaluation reward oscillates late in training (e.g.\ 36.2 at 180k vs.\ 14.0 at 190k steps), a known instability of DQN's
moving target; we therefore keep the best evaluation checkpoint.
% TODO [STEP 12]: convergence episode, sample efficiency, regret value
% TODO [STEP 10b]: epsilon-decay comparison (10% / 30% / 60%)
% TODO [STEP 9]: real-image (DeepCrack) result

% ---- add to references.bib ----
% @article{mnih2015human, title={Human-level control through deep reinforcement learning}, author={Mnih, Volodymyr and others},
%   journal={Nature}, volume={518}, pages={529--533}, year={2015}}
% @article{raffin2021sb3, title={Stable-Baselines3: Reliable Reinforcement Learning Implementations}, author={Raffin, Antonin and others},
%   journal={Journal of Machine Learning Research}, volume={22}, number={268}, pages={1--8}, year={2021}}
